\documentclass[11pt]{article}
% ===========================================================================
%  Shared preamble for the AP Calculus AB exam-preparation worksheets.
%  Input right after \documentclass. Files beginning with "_" are not compiled
%  on their own. Navy and gold, to match the course's lesson pages.
% ===========================================================================
\usepackage[letterpaper,top=2.3cm,bottom=2.2cm,left=1.7cm,right=1.7cm,headheight=28pt,footskip=30pt]{geometry}
\usepackage{amsmath,amssymb}
\usepackage[T1]{fontenc}
\usepackage[utf8]{inputenc}
\usepackage{lmodern}
\usepackage{xcolor}
\usepackage[most]{tcolorbox}
\usepackage{enumitem}
\usepackage{fancyhdr}
\usepackage{array}
\usepackage{tikz}
\usepackage{pgfplots}
\pgfplotsset{compat=1.18}

\definecolor{navy}{HTML}{161E54}
\definecolor{ink}{HTML}{27357E}
\definecolor{gold}{HTML}{B98A00}
\definecolor{goldbg}{HTML}{FFFDF3}
\definecolor{inkbg}{HTML}{F6F7FD}
\definecolor{rule}{HTML}{C9CFF0}
\definecolor{okgreen}{HTML}{1B7A44}
\definecolor{okbg}{HTML}{F1F8F3}

\pagestyle{fancy}
\fancyhf{}
\fancyhead[L]{\footnotesize Name: \underline{\hspace{4cm}}}
\fancyhead[C]{\textbf{\color{navy}Integration Academy}\\[-2pt]{\scriptsize Calculus AB \textperiodcentered{} AP\textsuperscript{\textregistered} Exam Preparation}}
\fancyfoot[L]{\tiny\color{black!55}\parbox[t]{0.78\textwidth}{AP\textsuperscript{\textregistered} is a registered trademark of the College Board, which was not involved in the production of, and does not endorse, this material. \textcopyright{} Integration Academy.}}
\fancyfoot[R]{\scriptsize\thepage}
\renewcommand{\headrulewidth}{1.2pt}
\renewcommand{\headrule}{\hbox to\headwidth{\color{gold}\leaders\hrule height \headrulewidth\hfill}}
\renewcommand{\footrulewidth}{0pt}
\setlength{\parindent}{0pt}
\setlength{\parskip}{4pt}

% Title block: \wstitle{code}{title}{unit line}
\newcommand{\wstitle}[3]{\begin{center}{\small\bfseries\color{gold}WORKSHEET #1}\\[3pt]{\LARGE\bfseries\color{navy}#2}\\[3pt]{\small\color{black!60}#3}\end{center}\vspace{2pt}}
% Section heading
\newcommand{\sech}[1]{\par\vspace{10pt}{\large\bfseries\color{navy}#1}\par\nobreak\vspace{2pt}{\color{rule}\hrule height0.8pt}\vspace{7pt}}
% A level tag printed beside a group of questions
\newcommand{\level}[1]{\par\vspace{6pt}{\footnotesize\bfseries\color{gold}\MakeUppercase{#1}}\par\nobreak\vspace{2pt}}

\newtcolorbox{ideabox}{colframe=ink,colback=inkbg,coltitle=white,colbacktitle=ink,fonttitle=\bfseries\small,title={Key ideas},arc=2pt,boxrule=0.8pt,left=8pt,right=8pt}
\newtcolorbox{example}[1]{colframe=rule,colback=white,coltitle=navy,colbacktitle=inkbg,fonttitle=\bfseries\small,title={#1},arc=2pt,boxrule=0.8pt,borderline west={3pt}{0pt}{ink},left=10pt,right=8pt}
\newtcolorbox{frbox}[1]{colframe=gold,colback=goldbg,coltitle=white,colbacktitle=gold,fonttitle=\bfseries\small,title={#1},arc=2pt,boxrule=0.8pt,breakable,left=8pt,right=8pt}
\newtcolorbox{solbox}{colframe=okgreen,colback=okbg,coltitle=white,colbacktitle=okgreen,fonttitle=\bfseries\small,title={Worked solutions},arc=2pt,boxrule=0.8pt,breakable,left=8pt,right=8pt}
\newtcolorbox{scorebox}[1]{colframe=navy,colback=white,coltitle=white,colbacktitle=navy,fonttitle=\bfseries\small,title={#1},arc=2pt,boxrule=0.8pt,breakable,left=8pt,right=8pt}

\newcommand{\soln}{\par\smallskip\textbf{\color{ink}Solution.}\ }
\newcommand{\result}[1]{\par\smallskip\textbf{#1}}
% blank working space of a given height
\newcommand{\work}[1]{\par\vspace{#1}}
% a point value in the margin of a scoring guide
\newcommand{\pts}[1]{\hfill{\footnotesize\bfseries\color{gold}[#1]}}

% Coordinate grid for graph questions: \begin{apgraph}{xmin}{xmax}{ymin}{ymax} ... \end{apgraph}
\newenvironment{apgraph}[4]{\begin{center}\begin{tikzpicture}\begin{axis}[
  axis lines=middle, xmin=#1, xmax=#2, ymin=#3, ymax=#4,
  grid=both, grid style={black!12}, xtick distance=1, ytick distance=1,
  ticklabel style={font=\scriptsize}, xlabel={$x$}, ylabel={$y$},
  width=10cm, height=7.2cm, samples=80, clip=false,
  every axis plot/.append style={line width=1.1pt, color=ink}]}{\end{axis}\end{tikzpicture}\end{center}}
\newcommand{\closeddot}[2]{\addplot[only marks, mark=*, mark size=2.2pt, color=ink] coordinates {(#1,#2)};}
\newcommand{\opendot}[2]{\addplot[only marks, mark=*, mark size=2.2pt, color=ink, mark options={fill=white}] coordinates {(#1,#2)};}

% Multiple-choice options: four across, or stacked when the options are long.
\newcommand{\opts}[4]{\par\vspace{3pt}\begin{tabular}{@{}p{0.235\linewidth}p{0.235\linewidth}p{0.235\linewidth}p{0.235\linewidth}@{}}(A)\ #1 & (B)\ #2 & (C)\ #3 & (D)\ #4\end{tabular}\par\vspace{9pt}}
\newcommand{\optsv}[4]{\par\vspace{3pt}\begin{tabular}{@{}l@{\ \ }p{0.9\linewidth}@{}}(A) & #1\\(B) & #2\\(C) & #3\\(D) & #4\end{tabular}\par\vspace{9pt}}

\fancyhead[R]{\footnotesize\color{black!60}Worksheet 5.7}
\begin{document}

\wstitle{5.7}{The Second Derivative Test}{Unit 5: Analysing Functions}

\begin{ideabox}
\begin{itemize}[leftmargin=*,itemsep=2pt,topsep=2pt]
  \item If $f'(c)=0$ and $f''(c)>0$, then $f$ has a relative \textbf{minimum} at $c$.
  \item If $f'(c)=0$ and $f''(c)<0$, then $f$ has a relative \textbf{maximum} at $c$.
  \item If $f'(c)=0$ and $f''(c)=0$, the test is \textbf{inconclusive}: use the first derivative test.
  \item If $c$ is the only critical point and $f''$ has one sign on the whole interval, the extremum is absolute.
\end{itemize}
\end{ideabox}

\sech{Worked examples}
\begin{example}{Example 1 \textperiodcentered{} Two critical points}
Classify the critical points of $f(x)=x^3-3x^2+2$.\soln $f'(x)=3x(x-2)$ and $f''(x)=6x-6$. $f''(0)=-6$ and $f''(2)=6$.\result{Relative maximum at $x=0$; relative minimum at $x=2$.}
\end{example}

\begin{example}{Example 2 \textperiodcentered{} When the test fails}
What does the test say about $f(x)=x^5$ at $x=0$?\soln $f'(0)=0$ and $f''(0)=0$: inconclusive. $f'(x)=5x^4$ is positive on both sides.\result{The test gives no information; in fact there is no extremum.}
\end{example}

\begin{example}{Example 3 \textperiodcentered{} An absolute minimum}
Find the absolute minimum of $f(x)=2x+\dfrac8x$ for $x>0$.\soln $f'(x)=2-\dfrac{8}{x^2}=0$ at $x=2$, the only critical point. $f''(x)=\dfrac{16}{x^3}>0$ for all $x>0$.\result{The absolute minimum is $f(2)=8$.}
\end{example}

\newpage
\sech{Practice}

\level{Core}
\begin{enumerate}[leftmargin=*,itemsep=2pt]
  \item Given $f'(3)=0$ and $f''(3)=2$, what does $f$ have at $x=3$?\work{2.3cm}
  \item Given $f'(1)=0$ and $f''(1)=-5$, what does $f$ have at $x=1$?\work{2.3cm}
  \item Use the second derivative test on $f(x)=x^2-8x$.\work{2.3cm}
  \item Given $f'(c)=0$ and $f''(c)=0$, what can be concluded?\work{2.3cm}
\end{enumerate}

\level{Exam level}
\begin{enumerate}[leftmargin=*,itemsep=2pt,start=5]
  \item Use the second derivative test to classify the critical points of $f(x)=x^3-48x$.\work{2.8cm}
  \item Find and classify the critical point of $f(x)=x^2-\ln x$ for $x>0$.\work{2.8cm}
  \item Use the second derivative test to classify the critical points of $f(x)=2\sin x+x$ on $(0,2\pi)$.\work{2.8cm}
  \item A solution of $\dfrac{dy}{dx}=x^2-y$ passes through the point $(2,4)$. Show that it has a relative minimum there.\work{2.8cm}
\end{enumerate}

\begin{frbox}{Question 9 \textperiodcentered{} Free-response practice}
Let $f(x)=\dfrac{x^4}{4}-2x^2$.
\begin{enumerate}[label=(\alph*),leftmargin=*,itemsep=2pt]
  \item Find the critical points of $f$.\work{2.6cm}
  \item Use the second derivative test to classify each critical point.\work{2.6cm}
  \item Find the relative maximum value and the relative minimum value of $f$.\work{2.6cm}
\end{enumerate}
\end{frbox}

\newpage
\begin{solbox}
\begin{enumerate}[leftmargin=*,itemsep=4pt]
  \item A relative minimum.
  \item A relative maximum.
  \item $f'(4)=0$ and $f''(4)=2>0$: relative minimum at $x=4$.
  \item Nothing: the test is inconclusive.
  \item $f'(x)=3x^2-48=0$ at $x=\pm4$; $f''(x)=6x$. Minimum at $x=4$, maximum at $x=-4$.
  \item $f'(x)=2x-\dfrac1x=0$ at $x=\dfrac{1}{\sqrt2}$. $f''(x)=2+\dfrac{1}{x^2}>0$: a relative (and absolute) minimum.
  \item $f'(x)=2\cos x+1=0$ at $\dfrac{2\pi}{3}$ and $\dfrac{4\pi}{3}$. $f''(x)=-2\sin x$: $-\sqrt3$ at $\dfrac{2\pi}{3}$ (maximum), $\sqrt3$ at $\dfrac{4\pi}{3}$ (minimum).
  \item At $(2,4)$, $\dfrac{dy}{dx}=4-4=0$, and $\dfrac{d^2y}{dx^2}=2x-\dfrac{dy}{dx}=4>0$.
  \item \begin{enumerate}[label=(\alph*),leftmargin=*,itemsep=3pt]
    \item $f'(x)=x^3-4x=x(x-2)(x+2)$: $x=0$ and $x=\pm2$.
    \item $f''(x)=3x^2-4$. $f''(0)=-4<0$: maximum. $f''(\pm2)=8>0$: minima.
    \item Maximum value $f(0)=0$; minimum value $f(\pm2)=4-8=-4$.
  \end{enumerate}
\end{enumerate}
\end{solbox}

\end{document}
